\subsection{Diversification history of the New World Monkeys}

BELLA analyses place the origin of platyrrhines around 40 Ma in a small-bodied ancestor weighing less than 500 g (category 0, \textit{Tiny}, in our discretized body-mass framework). This estimate is consistent with interpretations of the earliest fossil evidence for the clade \citep{bond2015eocene}.
Subsequently, the clade diversified into multiple larger-bodied lineages, some of which exceeded 3000 g (category 3, \textit{Large}) by the early Miocene (ca. 20 Ma; Fig. \ref{fig:platyrrhine}a).
This reconstruction corroborates previous inferences based on continuous models of trait evolution \citep{silvestro2019early}.

The evolution of platyrrhines is characterized by substantial variation in diversification rates, which our BELLA model attributes to heterogeneities through time and among body-mass categories.
Temporal effects were strongest for extinction, with rates peaking during the late Miocene (11–7 Ma), consistent with previous findings \citep{silvestro2019early}, with large-bodied lineages being more strongly affected.
This extinction peak is likely linked to the combined effects of global cooling following the Middle Miocene Climatic Optimum and the aridification of southern South America associated with Andean uplift and rain-shadow effects \citep{hoorn2022editorial}.
Together, these results point to a non-additive interaction between temporal and trait-dependent effects on extinction, a pattern that would be difficult to detect using alternative phylodynamic models.

The sustained decline in extinction rates, together with a moderate increase in speciation rates, leads to strongly positive net diversification toward the present, in agreement with previous findings based on simpler time-dependent phylodynamic models applied to extant-species phylogenies \citep{herrera2017primate}.
Speciation and extinction rates through time are also remarkably consistent with paleodiversity estimates inferred from an independently compiled dataset of fossil occurrences and derived using different methodologies and assumptions \citep{pino2025splendid}.
Notably, that analysis identified a rapid decline in predicted species richness during the late Miocene, followed by a rebound—dynamics that our results here attribute to a simultaneous increase in speciation and decrease in extinction.

While the credible intervals around the marginal rates through time remain wide (Fig. \ref{fig:platyrrhine:marginal}), they are substantially narrower than those obtained in a previous FBD analysis based on the same phylogenetic trees, which assumed a piecewise-constant model with independent rate parameters across time bins regularized by a Brownian process prior \citep{silvestro2019early}.
In particular, the upper bounds of the rates (here approximately 0.8 for speciation and 0.6 for extinction events per lineage per million years) are roughly tenfold lower than the highest values estimated previously.
This highlights BELLA’s ability to yield more tightly constrained uncertainty intervals than predictor-agnostic piecewise-constant models, a pattern also observed in simulation analyses (Figs. \ref{fig:epi-skyline:ribbons}, \ref{fig:fbd-no-traits:ribbons}).